% !TeX root = main.tex
\section{工业数据科学}

\subsection{简介}

理解数据分析能做到什么程度，常用一条从低到高的价值线：
\begin{enumerate}
    \item 描述性分析（Descriptive）：回答发生了什么，是最基础的观察和统计
    \item 诊断性分析（Diagnostic）：回答为什么会发生，需要解释原因
    \item 预测性分析（Predictive）：回答接下来会发生什么，需要对未来做预判
    \item 规范性分析（Prescriptive）：回答应该怎么做，直接给出行动建议
\end{enumerate}
越高的等级能带来的业务价值越大，技术难度也越高。描述这类海量、复杂数据的特征，通常用\textbf{3V}：
\begin{itemize}
    \item Volume（容量）
    \item Velocity（速度）
    \item Variety（多样性）
\end{itemize}


\subsection{CRISP-DM}

接下来我们用 \textbf{CRISP-DM（Cross-Industry Standard Process for Data Mining，跨行业数据挖掘标准流程）}来讲解一个喷油器的案例。这套方法论把一个完整的数据科学项目拆成六个阶段：
\begin{enumerate}
    \item \textbf{业务理解（Business Understanding）}：搞清楚项目到底要解决什么问题
    \item \textbf{数据理解（Data Understanding）}：搞清楚手头有什么数据、质量如何
    \item \textbf{数据准备（Data Preparation）}：把原始数据加工成能喂给模型的形式
    \item \textbf{建模（Modeling）}：选择合适的算法，训练模型
    \item \textbf{评估（Evaluation）}：判断模型到底好不好用
    \item \textbf{部署（Deployment）}：把模型真正用到实际生产决策中去
\end{enumerate}
这六个阶段不是一次性走完就结束的直线流程，而是一个可以循环迭代的闭环。评估阶段发现问题，往往要回头重新做数据准备甚至业务理解。

喷油器的生产分为部件加工和总装两大部分。总装线上要经历一连串工序，最后要经过一套完整的液压终检。这是一种 100\% 全检，每个产品要在多个检测点（P1、P2……）上分别测量喷射量、压力、回流量等多项特性，再和规格比对，最终打上合格标签。问题在于这种全检非常耗时，成了整条生产线的瓶颈工序。项目要解决的目标很明确：通过数据分析提前预测最终检测点的结果，从而减轻检测环节的压力。

拿到原始数据后往往不能直接拿去训练模型，通常要先完成几件事：
\begin{itemize}
    \item 数据清洗：识别并处理不可用的数值。
    \item 数据变换：不同的数据往往处在不同的量纲层级上。需要做升级或降级处理，让不同类型的数据能放在一起比较。
    \item 特征提取：每个检测点采集到的原始曲线通过特征提取转换成模型可以使用的特征。
\end{itemize}

数据准备好之后，就进入建模阶段。训练模型可以有多种机器学习的方法，例如线性回归、逻辑回归、k近邻、支持向量机等等，这些机器学习的方法会在第\ref{sec:aiml}章节中详细介绍。模型训练完了，还需要一套指标来判断它靠不靠谱。标准做法是用\textbf{混淆矩阵}框架，把预测结果和真实结果交叉对比，划分出四种情况：
\begin{itemize}
    \item \textbf{阳性（TP）}：预测不合格，实际也确实不合格
    \item \textbf{假阳性（Pseudofehlerrate, FP）}：预测不合格，但实际是合格的，伪错误。意味着模型"多疑"，把好产品误判成了坏产品
    \item \textbf{假阴性（Schlupfrate, FN）}：预测合格，但实际是不合格的，这类错误被称为漏检，是更危险的一种错误，意味着有缺陷产品被放过去了
    \item \textbf{真阴性（TN）}：预测合格，实际也确实合格
\end{itemize}
在这四个数字的基础上，可以算出一整套评估指标：
\begin{itemize}
    \item\textbf{精确率（Precision）}衡量被预测为合格的产品里，有多少是真的合格
    \item \textbf{召回率（Recall）}衡量所有真正不合格的产品里，有多少被成功抓出来了
    \item \textbf{特异度（Specificity）}衡量所有真正合格的产品里，有多少被正确放行了
    \item \textbf{准确率（Accuracy）}则是预测正确的比例，也就是 (TP+TN) / 全部样本数
    \item \textbf{伪错误率}：FP 占比
    \item \textbf{漏检率}：FN 占比
\end{itemize}
喷油器数据上跑出的结果整体准确率达到 96.15\%，但召回率只有 33.83\%，也就是说，虽然总体判断准确，但真正能抓出来的不合格品其实只占三分之一左右。这说明单看准确率会掩盖问题，必须结合多个指标一起看。

漏检率和伪错误率之间存在此消彼长的权衡。想把漏检率压得更低，往往就得接受更高的伪错误率，反之亦然。选哪个模型、定在哪个平衡点，最终要看业务上更能承受哪种代价。

